% 《物理实验报告模板.tex》的共同版式；所有实验模板均载入本文件。
\documentclass[UTF8,a4paper,10pt,fontset=fandol]{ctexart}
\usepackage[left=3.175cm,right=3.175cm,top=2.54cm,bottom=2.54cm]{geometry}
\usepackage{graphicx,tikz,pgfplots,needspace,adjustbox,booktabs,fancyhdr,lastpage,amsmath,amssymb,array,tabularx,longtable}
\pgfplotsset{compat=1.18}
\setlength{\parindent}{0pt}
\setlength{\parskip}{2pt}
\setlength{\headheight}{18pt}
\linespread{1}
\newcommand{\SetField}[2]{\expandafter\def\csname reportfield@#1\endcsname{#2}}
\makeatletter
\newcommand{\Field}[2][1.35cm]{%
  \ifcsname reportfield@#2\endcsname
    \expandafter\ifx\csname reportfield@#2\endcsname\@empty
      \underline{\makebox[#1]{\strut}}%
    \else\csname reportfield@#2\endcsname\fi
  \else\underline{\makebox[#1]{\strut}}\fi}
\makeatother
\newcommand{\MathField}[2][1.35cm]{%
  \ifcsname reportfield@math_#2\endcsname
    \csname reportfield@math_#2\endcsname
  \else\text{\Field[#1]{#2}}\fi}
\newcommand{\coverfield}[2]{%
  \noindent{\heiti\fontsize{18pt}{48pt}\selectfont
  \makebox[3.3cm][l]{#1}\underline{\makebox[8.8cm][c]{\Field[0pt]{#2}}}\par}}
\newcommand{\bindingline}[1][2.85]{%
  \begin{tikzpicture}[remember picture,overlay]
    \foreach \i in {0,...,38}{%
      \pgfmathsetmacro{\bindingy}{-#1-0.52*\i}%
      \ifnum\i=13\relax
        \node[inner sep=0pt] at ([xshift=1.05cm,yshift=\bindingy cm]current page.north west) {\songti\fontsize{10.5pt}{10.5pt}\selectfont 装};
      \else\ifnum\i=19\relax
        \node[inner sep=0pt] at ([xshift=1.05cm,yshift=\bindingy cm]current page.north west) {\songti\fontsize{10.5pt}{10.5pt}\selectfont 订};
      \else\ifnum\i=25\relax
        \node[inner sep=0pt] at ([xshift=1.05cm,yshift=\bindingy cm]current page.north west) {\songti\fontsize{10.5pt}{10.5pt}\selectfont 线};
      \else
        \pgfmathsetmacro{\bindingtop}{\bindingy+0.20}%
        \pgfmathsetmacro{\bindingbottom}{\bindingy-0.20}%
        \draw[line width=0.35pt,dash pattern=on 1.2pt off 1.8pt]
          ([xshift=1.05cm,yshift=\bindingtop cm]current page.north west)
          -- ([xshift=1.05cm,yshift=\bindingbottom cm]current page.north west);
      \fi\fi\fi}
  \end{tikzpicture}%
}
% 装订线印在每一页正文和原始数据页，续页也沿用母模板样式。
\newif\ifreportbinding
\AddToHook{shipout/foreground}{\ifreportbinding\bindingline\fi}
\fancypagestyle{report}{\fancyhf{}\fancyfoot[C]{第\thepage 页\quad 共\pageref{LastPage}页}\renewcommand{\headrulewidth}{0.4pt}}
\fancypagestyle{record}{\fancyhf{}\fancyhead[C]{原始数据记录单}\fancyfoot[C]{第\thepage 页\quad 共\pageref{LastPage}页}\renewcommand{\headrulewidth}{0.4pt}}
\newcommand{\ReportCover}[1]{%
  \thispagestyle{empty}\vspace*{1.0cm}
  \begin{center}
    \includegraphics[width=8.184cm]{tongji_logo.jpeg}\par
    {\bfseries\fontsize{18pt}{24pt}\selectfont TONGJI UNIVERSITY\par}
    \vspace{1.5cm}{\heiti\fontsize{36pt}{44pt}\selectfont 物理实验报告\par}
  \end{center}
  \vspace{1.5cm}
  \coverfield{实验名称}{experiment}
  \coverfield{学生姓名}{student}
  \coverfield{学\hspace{1em}号}{student_id}
  \coverfield{组\hspace{1em}号}{group}
  \coverfield{实验日期}{date}
  \clearpage\reportbindingtrue\pagestyle{report}
  \begin{center}{\heiti\fontsize{14pt}{40pt}\selectfont #1\par}\end{center}}
\newcommand{\ReportSection}[1]{\Needspace{4\baselineskip}\par\vspace{0.35\baselineskip}{\heiti #1}\par}
\newcommand{\ReportSubsection}[1]{\Needspace{3\baselineskip}\par\vspace{0.2\baselineskip}{\bfseries #1}\par}
\newcommand{\ReportRecord}{%
  \clearpage\thispagestyle{record}
  \vspace*{0.6cm}（请将教师签字的实验数据记录单完整拍照粘贴至此处。）
}

% ===== 数据输入区开始 =====
\SetField{experiment}{用波尔共振仪研究物体受迫振动}
% 后端以 \SetField{字段名}{值} 填入数据；缺省值保持空白。
% 母模板封面字段：\Field{experiment} \Field{student} \Field{student_id} \Field{group} \Field{date}
% ===== 数据输入区结束 =====
\begin{document}\fontsize{10.5pt}{20pt}\selectfont
\ReportCover{用波尔共振仪研究物体受迫振动}
\ReportSection{实验目的}
1. 测量摆轮的自由周期和阻尼系数。\par
2. 研究受迫振动的幅频、相频特性，观察共振现象。\par
3. 学习用频闪法测量相位差。
\ReportSection{实验原理}
摆轮在弹簧恢复力、阻尼力和周期外力作用下做受迫振动，其运动方程为
\[
\ddot\theta+2\beta\dot\theta+\omega_0^2\theta=m\cos\omega t.
\]
其中，$\theta$ 为角位移，$\beta=b/(2J)$ 为阻尼系数，$\omega_0=\sqrt{k/J}$ 为固有圆频率，$\omega=2\pi/T$ 为外力圆频率，$m=M_0/J$。$J$ 为转动惯量，$k$ 为弹性力矩系数，$b$ 为阻尼力矩系数，$M_0$ 为外力矩幅值。
稳定后，摆轮的振幅 $\theta_2$ 和相位差 $\varphi$ 满足
\[
\theta_2=\frac{m}{\sqrt{(\omega_0^2-\omega^2)^2+4\beta^2\omega^2}},\qquad
\tan\varphi=-\frac{2\beta\omega}{\omega_0^2-\omega^2}.
\]
外力频率接近固有频率时，振幅明显增大，出现共振。摆轮的振动落后于外力，因此上式中的 $\varphi$ 为负；本报告的表格和相频图统一用它的大小 $\phi=|\varphi|$，即仪器的正读数。
\ReportSection{实验仪器}
波尔共振仪、光电门、频闪相位读盘。
\ReportSection{实验内容}
1. 无电磁阻尼时，记录摆轮从约 $160^\circ$ 衰减至 $50^\circ$ 过程中的振幅及自由周期。\par
2. 选择阻尼档位，记录十个连续振幅及十周期总时间，重复测量。\par
3. 保持阻尼档位不变，调节外力周期；每次待振幅稳定后，记录周期、振幅和相位差。共振附近适当加密测量点。
\clearpage\ReportSection{数据记录与处理}
\ReportSubsection{自由振动}
下表给出振幅与自由周期 $T_0$ 的对应关系。计算某一振幅下的理论相位时，取该振幅对应的 $T_0$；没有相同振幅时，用相邻两点内插。
\begin{center}\small\renewcommand{\arraystretch}{1.35}
\begin{tabular}{|c|c|c|}\hline
序号 & 振幅/° & $T_0$/s\\\hline
1 & \Field[10mm]{free_angle_01} & \Field[10mm]{free_period_01}\\\hline
2 & \Field[10mm]{free_angle_02} & \Field[10mm]{free_period_02}\\\hline
3 & \Field[10mm]{free_angle_03} & \Field[10mm]{free_period_03}\\\hline
4 & \Field[10mm]{free_angle_04} & \Field[10mm]{free_period_04}\\\hline
5 & \Field[10mm]{free_angle_05} & \Field[10mm]{free_period_05}\\\hline
6 & \Field[10mm]{free_angle_06} & \Field[10mm]{free_period_06}\\\hline
7 & \Field[10mm]{free_angle_07} & \Field[10mm]{free_period_07}\\\hline
8 & \Field[10mm]{free_angle_08} & \Field[10mm]{free_period_08}\\\hline
9 & \Field[10mm]{free_angle_09} & \Field[10mm]{free_period_09}\\\hline
10 & \Field[10mm]{free_angle_10} & \Field[10mm]{free_period_10}\\\hline
\end{tabular}
\end{center}
\ReportSubsection{阻尼系数}
阻尼档位：$\Field{damping_setting}$。阻尼振动单周期 $T_d=\Field{damped_period_s}$ s，由十周期总时间除以十得到。相隔五个周期的振幅按逐差法计算：
\[
\beta=\frac{1}{5T_d}\,\overline{\ln\frac{\theta_i}{\theta_{i+5}}},\qquad i=1,2,\ldots,5.
\]
先求五个振幅比的自然对数，再取平均，最后除以 $5T_d$。多次测量分别计算后取平均。
\begin{center}\small\renewcommand{\arraystretch}{1.35}
\begin{tabular}{|c|c|c|c|c|}\hline
次数 & $\theta_i$/° & 次数 & $\theta_{i+5}$/° & $\ln(\theta_i/\theta_{i+5})$\\\hline
1 & \Field[10mm]{damped_peak_01} & 6 & \Field[10mm]{damped_peak_06} & \Field[10mm]{damped_log_ratio_01}\\\hline
2 & \Field[10mm]{damped_peak_02} & 7 & \Field[10mm]{damped_peak_07} & \Field[10mm]{damped_log_ratio_02}\\\hline
3 & \Field[10mm]{damped_peak_03} & 8 & \Field[10mm]{damped_peak_08} & \Field[10mm]{damped_log_ratio_03}\\\hline
4 & \Field[10mm]{damped_peak_04} & 9 & \Field[10mm]{damped_peak_09} & \Field[10mm]{damped_log_ratio_04}\\\hline
5 & \Field[10mm]{damped_peak_05} & 10 & \Field[10mm]{damped_peak_10} & \Field[10mm]{damped_log_ratio_05}\\\hline
\end{tabular}
\end{center}
\clearpage\ReportSubsection{幅频、相频特性}
$T$ 为外力周期，$\theta$ 为稳定振幅。最大振幅记为 $\theta_r$，对应周期记为 $T_r$；若最大振幅有多个点，取这些点周期的平均值。本次 $\theta_r=\Field{peak_amplitude_deg}$°，$T_r=\Field{resonance_period_s}$ s。
\[
\frac{\omega}{\omega_r}=\frac{T_r}{T},\qquad
\frac{\theta^2}{\theta_r^2}=\left(\frac{\theta}{\theta_r}\right)^2.
\]
相位计算使用该振幅对应的自由周期 $T_0$ 和测得的 $\beta$：
\[
\varphi=\tan^{-1}\!\left[\frac{-\beta T_0^2T}{\pi(T^2-T_0^2)}\right].
\]
按角度计算；若反正切结果为正，则减去 $180^\circ$；$T=T_0$ 时取 $-90^\circ$。下表“计算值”取其绝对值，与仪器正读数比较。若振幅超出自由振动表范围，标为“需补测”。
\begin{center}\small\renewcommand{\arraystretch}{1.35}
\begin{adjustbox}{max width=\textwidth}
\begin{tabular}{|c|c|c|c|c|c|}\hline
强迫力矩周期 (s) & \shortstack{相位差 $\phi$ ($^\circ$)\\读取值} & \shortstack{$\theta$ ($^\circ$)\\测量值} & $\omega/\omega_r$ & $\theta^2/\theta_r^2$ & \shortstack{相位差 $\phi$ ($^\circ$)\\计算值}\\\hline
\Field[10mm]{table_drive_period_01} & \Field[10mm]{table_phase_reading_01} & \Field[10mm]{amplitude_01} & \Field[10mm]{table_frequency_ratio_01} & \Field[10mm]{table_amplitude_ratio_01} & \Field[10mm]{table_phase_calculated_01}\\\hline
\Field[10mm]{table_drive_period_02} & \Field[10mm]{table_phase_reading_02} & \Field[10mm]{amplitude_02} & \Field[10mm]{table_frequency_ratio_02} & \Field[10mm]{table_amplitude_ratio_02} & \Field[10mm]{table_phase_calculated_02}\\\hline
\Field[10mm]{table_drive_period_03} & \Field[10mm]{table_phase_reading_03} & \Field[10mm]{amplitude_03} & \Field[10mm]{table_frequency_ratio_03} & \Field[10mm]{table_amplitude_ratio_03} & \Field[10mm]{table_phase_calculated_03}\\\hline
\Field[10mm]{table_drive_period_04} & \Field[10mm]{table_phase_reading_04} & \Field[10mm]{amplitude_04} & \Field[10mm]{table_frequency_ratio_04} & \Field[10mm]{table_amplitude_ratio_04} & \Field[10mm]{table_phase_calculated_04}\\\hline
\Field[10mm]{table_drive_period_05} & \Field[10mm]{table_phase_reading_05} & \Field[10mm]{amplitude_05} & \Field[10mm]{table_frequency_ratio_05} & \Field[10mm]{table_amplitude_ratio_05} & \Field[10mm]{table_phase_calculated_05}\\\hline
\Field[10mm]{table_drive_period_06} & \Field[10mm]{table_phase_reading_06} & \Field[10mm]{amplitude_06} & \Field[10mm]{table_frequency_ratio_06} & \Field[10mm]{table_amplitude_ratio_06} & \Field[10mm]{table_phase_calculated_06}\\\hline
\Field[10mm]{table_drive_period_07} & \Field[10mm]{table_phase_reading_07} & \Field[10mm]{amplitude_07} & \Field[10mm]{table_frequency_ratio_07} & \Field[10mm]{table_amplitude_ratio_07} & \Field[10mm]{table_phase_calculated_07}\\\hline
\Field[10mm]{table_drive_period_08} & \Field[10mm]{table_phase_reading_08} & \Field[10mm]{amplitude_08} & \Field[10mm]{table_frequency_ratio_08} & \Field[10mm]{table_amplitude_ratio_08} & \Field[10mm]{table_phase_calculated_08}\\\hline
\Field[10mm]{table_drive_period_09} & \Field[10mm]{table_phase_reading_09} & \Field[10mm]{amplitude_09} & \Field[10mm]{table_frequency_ratio_09} & \Field[10mm]{table_amplitude_ratio_09} & \Field[10mm]{table_phase_calculated_09}\\\hline
\Field[10mm]{table_drive_period_10} & \Field[10mm]{table_phase_reading_10} & \Field[10mm]{amplitude_10} & \Field[10mm]{table_frequency_ratio_10} & \Field[10mm]{table_amplitude_ratio_10} & \Field[10mm]{table_phase_calculated_10}\\\hline
\Field[10mm]{table_drive_period_11} & \Field[10mm]{table_phase_reading_11} & \Field[10mm]{amplitude_11} & \Field[10mm]{table_frequency_ratio_11} & \Field[10mm]{table_amplitude_ratio_11} & \Field[10mm]{table_phase_calculated_11}\\\hline
\Field[10mm]{table_drive_period_12} & \Field[10mm]{table_phase_reading_12} & \Field[10mm]{amplitude_12} & \Field[10mm]{table_frequency_ratio_12} & \Field[10mm]{table_amplitude_ratio_12} & \Field[10mm]{table_phase_calculated_12}\\\hline
\Field[10mm]{table_drive_period_13} & \Field[10mm]{table_phase_reading_13} & \Field[10mm]{amplitude_13} & \Field[10mm]{table_frequency_ratio_13} & \Field[10mm]{table_amplitude_ratio_13} & \Field[10mm]{table_phase_calculated_13}\\\hline
\Field[10mm]{table_drive_period_14} & \Field[10mm]{table_phase_reading_14} & \Field[10mm]{amplitude_14} & \Field[10mm]{table_frequency_ratio_14} & \Field[10mm]{table_amplitude_ratio_14} & \Field[10mm]{table_phase_calculated_14}\\\hline
\Field[10mm]{table_drive_period_15} & \Field[10mm]{table_phase_reading_15} & \Field[10mm]{amplitude_15} & \Field[10mm]{table_frequency_ratio_15} & \Field[10mm]{table_amplitude_ratio_15} & \Field[10mm]{table_phase_calculated_15}\\\hline
\end{tabular}
\end{adjustbox}
\end{center}
\par\clearpage
\ReportSection{实验曲线}
两图均以 $\omega/\omega_r$ 为横轴，相位差取正值，与主要结果表一致。
\begin{center}\begin{tikzpicture}
\begin{axis}[width=.82\textwidth,height=6.2cm,xlabel={$\omega/\omega_r$},ylabel={$(\theta/\theta_r)^2$},grid=both]
\addplot+[mark=*,blue] coordinates {\PlotCoordinates{normalized_amplitude_graph_points}};
\end{axis}\end{tikzpicture}\par
\begin{tikzpicture}
\begin{axis}[width=.82\textwidth,height=6.2cm,xlabel={$\omega/\omega_r$},ylabel={相位差 $\phi$/$^\circ$},grid=both]
\addplot+[mark=square*,orange] coordinates {\PlotCoordinates{phase_magnitude_graph_points}};
\end{axis}\end{tikzpicture}\end{center}
\ReportSection{讨论与分析}
自由周期随振幅略有变化，计算理论相位时应按振幅查表，避免统一取平均周期造成误差。相位读数还受稳态等待时间、读盘零点和读数判断影响。共振附近变化较快，测量点过少容易错过最大振幅。
\ReportSection{实验结论}
逐差法测得阻尼系数 $\beta=\MathField{damping_beta_s_inv}\,\mathrm{s^{-1}}$。最大振幅为 $\MathField{peak_amplitude_deg}^{\circ}$，对应外力周期约 $T_r=\MathField{resonance_period_s}\,\mathrm{s}$。幅频、相频曲线反映了摆轮在共振附近振幅增大、相位差变化较快的特点。

\ifreportimages\ReportRecord\fi
% ===== 原始数据图片插入区 =====
\end{document}
